\documentclass[11pt]{article}

\usepackage[utf8]{inputenc}
\usepackage[T1]{fontenc}
\usepackage{lmodern}
\usepackage{geometry}
\usepackage{amsmath,amssymb,amsfonts,amsthm,mathtools}
\usepackage{bm}
\usepackage{enumitem}
\usepackage{microtype}
\usepackage{hyperref}
\usepackage{titlesec}
\usepackage{booktabs}
\usepackage{array}
\usepackage{setspace}
\usepackage{mathrsfs}
\usepackage{physics}

\geometry{margin=1in}
\hypersetup{colorlinks=true, linkcolor=black, urlcolor=blue, citecolor=black}
\setlist[enumerate]{topsep=0.4em,itemsep=0.35em}
\setlist[itemize]{topsep=0.3em,itemsep=0.25em}
\titleformat{\section}{\large\bfseries}{\thesection.}{0.5em}{}
\titleformat{\subsection}{\normalsize\bfseries}{\thesubsection.}{0.5em}{}
\setstretch{1.06}

\newcommand{\Aether}{\AE{}ther}
\newcommand{\Aflow}{\AE{}ther-flow}
\newcommand{\TheoryName}{\Aflow{} Interpretation of Relativity}
\newcommand{\TheoryProgram}{\Aether{} / \Aflow{} framework}
\newcommand{\CompletedMetric}{g^{\sharp}}

\newcommand{\KernelPackage}{\mathfrak F^{\mathrm{presel,ker}}}
\newcommand{\PreSectionSelectorPackage}{\mathfrak E^{\mathrm{presel,sel}}}
\newcommand{\PreSelectorMinSectionCore}{\mathfrak D^{\mathrm{presel,sec}}}
\newcommand{\YoctoSelectorPackage}{\mathfrak C^{\mathrm{yoc,sel}}}
\newcommand{\YoctoImageCore}{\mathfrak C^{\mathrm{yoc,img}}}
\newcommand{\YoctoMinSectionCore}{\mathfrak C^{\mathrm{yoc,sec}}}
\newcommand{\ZeptoImageCore}{\mathfrak B^{\mathrm{zep,img}}}
\newcommand{\ZeptoMinSectionCore}{\mathfrak B^{\mathrm{zep,sec}}}
\newcommand{\AttoSectionPackage}{\mathfrak A^{\mathrm{attosec}}}
\newcommand{\FemtoSectionPackage}{\mathfrak Z^{\mathrm{femtosec}}}
\newcommand{\PicoSectionPackage}{\mathfrak Y^{\mathrm{picosec}}}
\newcommand{\PicoMinSectionCore}{\mathfrak Y^{\mathrm{pic,sec}}}
\newcommand{\NanoSectionPackage}{\mathfrak X^{\mathrm{nanosec}}}
\newcommand{\MicroSectionPackage}{\mathfrak W^{\mathrm{microsec}}}
\newcommand{\SubSectionPackage}{\mathfrak V^{\mathrm{subsec}}}

\newtheorem{definition}{Definition}
\newtheorem{proposition}{Proposition}
\newtheorem{remark}{Remark}
\newtheorem{corollary}{Corollary}
\newtheorem{theorem}{Theorem}

\title{\textbf{The \TheoryName{}}\\[0.4em]
\large Primitive-Reservoir Observer-Reduction\\
\large Selector-Dynamics-Generating Substrate\\
\large Pre-Selector Kernel Dynamics\\
\large Next Benchmark Criterion}
\author{}
\date{}

\begin{document}

\maketitle

\begin{abstract}
The current active-primary theorem already removed the current pre-selector-minimizer-section-generating substrate selector dynamics package \(\PreSectionSelectorPackage\) as the live benchmark-facing datum by deriving it canonically from deeper selector-dynamics-generating substrate pre-selector kernel dynamics \(\KernelPackage\). This renews the live frontier. Another same-output deeper-origin relay that merely preserves \(\KernelPackage\) and therefore merely reproduces the same current selector-dynamics package \(\PreSectionSelectorPackage\), the same current observer-relevant pre-selector minimizer-section core \(\PreSelectorMinSectionCore\), the same current selector package \(\YoctoSelectorPackage\), the same current observer-relevant yocto-section minimizer-section core \(\YoctoMinSectionCore\), the same current observer-relevant yocto-section minimizer-image core \(\YoctoImageCore\), and the same downstream observer-relevant zepto-section, atto-section, femto-section, pico-section, nano-section, micro-section, and sub-section consequences no longer counts as the honest default next benchmark-facing manuscript. Neither does a re-expansion back to the current selector-dynamics package, the larger pre-selector package, or the older yocto-section presentations, and neither does a quantum branch that does not do that same benchmark-facing work.

This note records the routing consequence of that theorem. The next honest benchmark-facing move must either derive or justify the selector-dynamics-generating substrate pre-selector kernel dynamics \(\KernelPackage\) from still more primitive substrate structure in a way that changes benchmark-facing status, or prove a genuinely smaller theorem-level collapse strictly inside \(\KernelPackage\). The benchmark boundary remains fixed throughout. The one operative metric is still exactly \(\CompletedMetric\), the ordinary matter route is unchanged, there is no second operative metric, the low-energy matter law is not enlarged, and no stronger promotion language is licensed. The positive result is therefore routing discipline at the current pre-selector kernel-dynamics frontier.
\end{abstract}

\section{Introduction}

The active derivational program of \TheoryName{} is no longer at the current selector-dynamics package as its default stopping point. The current active-primary theorem already showed that the pre-selector-minimizer-section-generating substrate selector dynamics package \(\PreSectionSelectorPackage\) is not the minimal benchmark-facing datum at current scope, because that package is now generated canonically by the deeper selector-dynamics-generating substrate pre-selector kernel dynamics \(\KernelPackage\). That theorem therefore spent the selector-dynamics presentation as the live burden below the current pre-selector selector-dynamics criterion.

The present note records the routing consequence of that gain. It does not reopen the already-spent selector-dynamics package as if it were still the live unexplained datum, and it does not claim a completed first-principles substrate derivation of gravity. It records only which kind of next manuscript would now count as an honest benchmark-facing gain once the live burden has moved one level lower to the deeper kernel dynamics themselves.

\paragraph{Framework and claim boundary.}
\TheoryProgram{} denotes the broader \Aether{} / \Aflow{} framework built on the claims that the \Aether{} is the underlying four-dimensional substrate of reality, the \Aflow{} is its intrinsic ordered motion, observed three-dimensional space is the local experiential slice of that deeper substrate, S-time is the experienced order of change arising from matter, light, and the \Aflow{}, and observed expansion is the three-dimensional appearance of deeper four-dimensional ordered motion. Within that framework, \TheoryName{} denotes the exact-closure theory in which the effective gravitational dynamics are adopted to be exactly Einsteinian with universal matter coupling. In the active sequence, \emph{adoption} means use of that established relativistic dynamics without claiming substrate derivation, while \emph{derivation} is reserved for a first-principles recovery from explicit substrate variables.

The present note stays strictly inside that boundary. It records only which kind of next manuscript would now count as an honest benchmark-facing gain at the current pre-selector kernel-dynamics frontier.

\section{Fixed Inputs}

\begin{definition}[Pre-selector kernel-dynamics frontier inputs]
The criterion recorded here uses the following fixed inputs.
\begin{enumerate}
    \item The benchmark exact-closure package, which fixes the public one-metric GR target of the repository.
    \item The benchmark gatekeeping layer, including the recorded safeguard that blocks same-output deeper-origin relays from counting as new front-line progress once the live benchmark-relevant datum is unchanged.
    \item The earlier selector-dynamics routing criterion, which already recorded that the honest next move below the current selector-dynamics frontier had to derive or justify \(\PreSectionSelectorPackage\) or collapse it further.
    \item The current active-primary theorem, which proves that the current selector-dynamics package \(\PreSectionSelectorPackage\) is itself generated canonically by the deeper selector-dynamics-generating substrate pre-selector kernel dynamics \(\KernelPackage\).
\end{enumerate}
\end{definition}

\begin{remark}
The present note does not replace those manuscripts. It records the routing consequence of reading them together under the active safeguard against recursive depth drift.
\end{remark}

\section{Why the Current Burden Now Sits at Pre-Selector Kernel-Dynamics Scope}

\begin{proposition}[The live burden is renewed at pre-selector kernel-dynamics scope]
At the current stage of the primitive-reservoir observer-reduction line, the current active-primary theorem renews the live benchmark-facing burden at the level of the selector-dynamics-generating substrate pre-selector kernel dynamics \(\KernelPackage\).
\end{proposition}

\begin{proof}
That theorem changes benchmark-facing status because it removes the current selector-dynamics package \(\PreSectionSelectorPackage\) as the live unexplained stopping point at current scope. Once \(\PreSectionSelectorPackage\) is shown to arise canonically from the deeper kernel dynamics \(\KernelPackage\), the live burden no longer sits on the selector-dynamics package itself. It sits on the deeper package that generates it. The earlier selector-dynamics criterion remains preserved main-line history and routing handoff, but it no longer defines the live frontier once the current selector-dynamics package itself has been benchmark-facingly derived from deeper kernel dynamics.
\end{proof}

\begin{definition}[Benchmark-neutral deeper relay at current pre-selector kernel-dynamics scope]
A \emph{benchmark-neutral deeper relay at current pre-selector kernel-dynamics scope} is a manuscript that preserves the same benchmark one-metric output, the same ordinary matter route, the same selector-dynamics-generating substrate pre-selector kernel dynamics \(\KernelPackage\), the same current selector-dynamics package \(\PreSectionSelectorPackage\), the same current observer-relevant pre-selector minimizer-section core \(\PreSelectorMinSectionCore\), the same current selector package \(\YoctoSelectorPackage\), the same current observer-relevant yocto-section minimizer-section core \(\YoctoMinSectionCore\), the same current observer-relevant yocto-section minimizer-image core \(\YoctoImageCore\), the same current observer-relevant zepto-section minimizer-section core \(\ZeptoMinSectionCore\), the same current observer-relevant zepto-section minimizer-image core \(\ZeptoImageCore\), the same current atto-section package \(\AttoSectionPackage\), the same current femto-section package \(\FemtoSectionPackage\), the same current pico-section package \(\PicoSectionPackage\), the same current pico-section minimizer-section core \(\PicoMinSectionCore\), the same current nano-section package \(\NanoSectionPackage\), the same current micro-section package \(\MicroSectionPackage\), and the same current sub-section package \(\SubSectionPackage\) while merely relocating the unexplained substrate layer deeper, repackaging the same current deeper kernel package, or re-expanding the presentation back to the current selector-dynamics package, the larger pre-selector package, or older yocto-section presentations without a new compression, necessity, uniqueness, or comparably sharp routing gain.
\end{definition}

\section{The Current Pre-Selector Kernel-Dynamics Criterion}

\begin{definition}[Admissible next benchmark-facing manuscript at current pre-selector kernel-dynamics scope]
At the current pre-selector kernel-dynamics frontier, a manuscript counts as an admissible next benchmark-facing move only if it does at least one of the following:
\begin{enumerate}
    \item derives or justifies the selector-dynamics-generating substrate pre-selector kernel dynamics \(\KernelPackage\) from still more primitive substrate structure in a way that does more than merely relocate the unexplained layer deeper;
    \item proves a genuinely smaller theorem-level collapse strictly inside \(\KernelPackage\);
    \item does either of the previous two while preserving the same benchmark one-metric package, the same ordinary matter route, and the same claim boundary.
\end{enumerate}
\end{definition}

\begin{theorem}[Next honest benchmark-facing task at current pre-selector kernel-dynamics scope]
The next honest benchmark-facing manuscript below the current pre-selector kernel-dynamics frontier must either derive or justify the selector-dynamics-generating substrate pre-selector kernel dynamics \(\KernelPackage\) from still more primitive substrate structure in a benchmark-relevant way, or else prove a genuinely smaller theorem-level collapse strictly inside that deeper package. A manuscript that only repeats the same current kernel dynamics through another deeper relay, re-expands the discussion back to the current selector-dynamics package, the larger pre-selector package, or older yocto-section presentations without such a new gain, or invokes a quantum branch without doing the same benchmark-facing work does not satisfy the criterion.
\end{theorem}

\begin{proof}
By Proposition 1, the live burden already lies at the deeper kernel dynamics \(\KernelPackage\) rather than at the current selector-dynamics package \(\PreSectionSelectorPackage\). Therefore a subsequent manuscript must improve on that status rather than merely accompany it. A deeper-origin manuscript can do so only if it changes benchmark-facing status by more than depth alone. If it simply reproduces the same current kernel dynamics and their same downstream outputs, the routing safeguard classifies it as benchmark neutral. The remaining honest alternatives are therefore exactly those stated in the definition: either a more primitive derivation or justification of the current kernel dynamics themselves, or a sharper internal collapse strictly inside that deeper package.
\end{proof}

\begin{remark}
At present scope, the default next target is the first route unless a genuinely smaller observer-relevant theorem datum is isolated inside \(\KernelPackage\). No such sharper internal datum is yet fixed on the active line.
\end{remark}

\section{What the Next Move Should Not Be}

\begin{proposition}[Screened next moves at the current pre-selector kernel-dynamics frontier]
At the current pre-selector kernel-dynamics frontier, none of the following is the default next benchmark-facing move:
\begin{enumerate}
    \item another same-output deeper-origin relay that merely preserves the current kernel dynamics \(\KernelPackage\) and therefore merely reproduces the same current selector-dynamics package \(\PreSectionSelectorPackage\), the same current observer-relevant pre-selector minimizer-section core \(\PreSelectorMinSectionCore\), the same current selector package \(\YoctoSelectorPackage\), the same current yocto-section minimizer-section core \(\YoctoMinSectionCore\), the same current yocto-section minimizer-image core \(\YoctoImageCore\), the same current observer-relevant zepto-section minimizer-section core \(\ZeptoMinSectionCore\), the same current observer-relevant zepto-section minimizer-image core \(\ZeptoImageCore\), the same current atto-section package \(\AttoSectionPackage\), the same current femto-section package \(\FemtoSectionPackage\), the same current pico-section package \(\PicoSectionPackage\), the same current pico-section minimizer-section core \(\PicoMinSectionCore\), and the same downstream nano-section, micro-section, and sub-section consequences;
    \item another re-expansion back to the current selector-dynamics package, the full pre-selector package, or the older yocto-section presentations as if those larger presentations were again independently live burdens;
    \item a quantum branch that does not derive or justify the same deeper kernel dynamics, collapse them further, or close a benchmark derivation gate in a genuinely new way;
    \item a manuscript that introduces a second operative metric, enlarges the ordinary matter law, or uses stronger promotion language.
\end{enumerate}
\end{proposition}

\begin{proof}
Item (1) fails the preceding criterion because it leaves the renewed live burden unchanged. Item (2) likewise fails because it re-expands to larger already-spent presentations rather than removing remaining freedom in \(\KernelPackage\) or proving a smaller internal datum there. Item (3) does not directly address the live burden at current scope unless it performs the same benchmark-facing work named above. Item (4) violates the fixed claim boundary of the benchmark package and therefore cannot count as a disciplined next step on the active line. The benchmark package remains the exact one-metric GR package already fixed by adoption, so any admissible next manuscript must preserve that boundary.
\end{proof}

\begin{remark}
This screening statement is not a denial that those deeper or alternative manuscripts may matter strategically. It is a routing statement: they are not the default way to spend the next benchmark-facing move from the current pre-selector kernel-dynamics frontier unless they do the same benchmark-facing work named above.
\end{remark}

\section{Conclusion}

The positive result of this note is routing discipline at the current pre-selector kernel-dynamics frontier. The current active-primary theorem remains the live benchmark-facing gain because it already removes the current selector-dynamics package as the unexplained stopping point on the active line. The earlier selector-dynamics criterion remains preserved main-line history, but under the recursive-depth safeguard it is no longer itself the default current frontier.

The scope remains narrow and honest. The benchmark package is unchanged: the one operative metric is still exactly \(\CompletedMetric\), the ordinary matter route is unchanged, there is no second operative metric, and the low-energy matter law is not enlarged. Nothing here upgrades adoption to derivation or licenses stronger promotion language.

The next open burden is therefore clear. The next benchmark-facing manuscript should derive or justify the selector-dynamics-generating substrate pre-selector kernel dynamics from still more primitive substrate structure in a way that changes benchmark-facing status, unless the sharper honest gain available instead is a genuinely smaller theorem-level collapse strictly inside that deeper package. Another same-output deeper relay that merely preserves those kernel dynamics, a re-expansion back to the current selector-dynamics package, the larger pre-selector package, or older yocto-section presentations, or a quantum branch that does not do that same benchmark-facing work is not the clean next manuscript target at current scope.

\end{document}
